\chapter{Matriks}

Operasi matriks seperti penjumlahan dan perkalian
bersifat mekanis sehingga sangat cocok dikerjakan oleh
komputer.



%========================================
\section{Membentuk matriks}
Definisikan matriks dengan fungsi milik \Sage, yaitu \inlinecode{matrix}.
Ingat dari bagian tentang Metode Gauss bahwa fungsi ini menerima dua masukan.
Masukan pertama adalah himpunan bilangan yang menentukan skalar apa saja
yang dapat menjadi entri, sedangkan masukan kedua adalah entri-entri tersebut,
yang ditulis sebagai daftar baris.
Di sini, untuk masukan pertama kita akan menggunakan \Sagecmd{RDF} bagi
entri bilangan titik mengambang,
atau \Sagecmd{CDF} bagi entri bilangan kompleks $a+bi$ dengan
$a$ dan~$b$ berupa bilangan titik mengambang,
atau kita dapat menggunakan bilangan rasional, \Sagecmd{QQ}.
\begin{sagecommandline}
sage: A = matrix(RDF, [[1, 2], [3, 4]])
sage: A
sage: i = CDF(i)
sage: A = matrix(CDF, [[1+2*i, 3+4*i], [5+6*i, 7+8*i]])
sage: A
sage: A = matrix(QQ, [[1, 2], [3, 4]])
sage: A                           
\end{sagecommandline}
Simbol bawaan \Sage{} untuk akar kuadrat dari $-1$ adalah $i$.
Karena~$i$ digunakan untuk banyak hal dalam pemrograman,
sebelum menggunakannya untuk bilangan kompleks Anda perlu
mengaturnya ulang dengan \inlinecode{CDF(i)}.

Dalam bab ini, kecuali ada alasan untuk memilih yang lain,
kita akan menggunakan bilangan rasional, \inlinecode{QQ}, sebagai skalar,
karena bilangan rasional lebih mudah dibaca\Dash $1$ lebih mudah daripada
$1.0$\Dash dan karena matriks-matriks dalam buku biasanya mempunyai entri rasional.

Konstruktor \inlinecode{matrix} memungkinkan Anda menentukan banyaknya
baris dan kolom.
\begin{sagecommandline}
sage: B = matrix(QQ, 2, 3, [[1, 1, 1], [2, 2, 2]])  
sage: B
\end{sagecommandline}
Salah satu alasan untuk melakukannya adalah memeriksa masukan Anda.
Di sini ukuran yang ditentukan tidak cocok dengan entri karena
matriks yang diberikan hanya mempunyai dua baris.
\begin{lstlisting}[style=python]
sage: B = matrix(QQ, 3, 3, [[1, 1, 1], [2, 2, 2]])  
\end{lstlisting}
Pesan galat \Sage{} berbunyi
\inlinecode{Number of rows does not match up with specified number}.

Keadaan lain ketika Anda perlu mencantumkan
banyaknya baris dan kolom adalah saat memakai pintasan berikut
untuk memperoleh matriks identitas, yakni dengan menaruh $1$
di tempat daftar baris.
\begin{sagecommandline}
sage: I = matrix(QQ, 3, 3, 1)                     
sage: I
\end{sagecommandline}
Pintasan serupa menghasilkan matriks nol.
\begin{sagecommandline}
sage: Z = matrix(QQ, 2, 2, 0)
sage: Z
\end{sagecommandline}
Perbedaan antara pintasan ini dan pintasan sebelumnya adalah bahwa
\inlinecode{matrix(QQ, 3, 2, 1)} menghasilkan galat karena
matriks identitas harus berbentuk persegi.
\Sage{} mempunyai perintah untuk membuat matriks identitas
yang tidak mungkin menimbulkan galat ini.
\begin{sagecommandline}
sage: I = identity_matrix(3)
sage: I
\end{sagecommandline}

\Sage{} mempunyai banyak metode lain untuk matriks.
Misalnya, Anda dapat mentransposisi baris menjadi kolom atau menguji apakah
suatu matriks simetris, yaitu tidak berubah oleh transposisi.
\begin{sagecommandline}
sage: A = matrix(QQ, [[1, 2], [3, 4]])
sage: A.transpose()
sage: A.is_symmetric()
\end{sagecommandline}
Contoh lainnya, Anda dapat membuat matriks dengan memberikan pola
bagi entri-entrinya.
Dalam matriks yang terbentuk di sini, entri $a_{i,j}$ sama dengan jumlah
$i+j$.
\begin{sagecommandline}
sage: A = matrix(QQ, 2, 3, lambda x, y: x+y)
sage: A
\end{sagecommandline}
Satu contoh terakhir:
\Sage{} akan membuat matriks dengan entri acak, di sini berupa
bilangan titik mengambang dari 0 sampai 1.
\begin{sagecommandline}
sage: random_matrix(RDF, 3, min=0, max=1)
\end{sagecommandline}
(Perhatikan \inlinecode{RDF}.
Pilihan ini lebih sesuai bagi kita karena \Sage{} menyediakan
\inlinecode{random_matrix} yang lebih langsung digunakan dengan entri titik mengambang daripada
dengan entri rasional.
Rinciannya tersedia dalam referensi \Sage.)




%========================================
\section{Kombinasi linear}
Penjumlahan dan pengurangan matriks merupakan operasi yang wajar.
\begin{sagecommandline}
sage: A = matrix(QQ, [[1, 2], [3, 4]])
sage: B = matrix(QQ, [[1, 1], [2, -2]])
sage: A+B
sage: A-B
\end{sagecommandline}

\Sage{} mengetahui bahwa penjumlahan matriks dengan ukuran berbeda tidak terdefinisi.
\begin{lstlisting}
sage: A = matrix(QQ, [[1, 2], [3, 4]])
sage: C = matrix(QQ, [[0, 0, 2], [3, 2, 1]])
sage: A+C
\end{lstlisting}
Anda memperoleh galat yang baris terakhirnya memuat
\inlinecode{unsupported operand parent(s) for +}.
Singkatnya, \inlinecode{+} tidak terdefinisi antara sebuah matriks
$\nbyn{2}$ dan sebuah matriks $\nbym{2}{3}$.

Perkalian skalar juga merupakan operasi yang wajar,
sehingga Anda dapat membentuk kombinasi linear.
\begin{sagecommandline}[d,0,2]
sage: A = matrix(QQ, [[1, 2], [3, 4]])
sage: B = matrix(QQ, [[1, 1], [2, -2]])
sage: 3*A
sage: 3*A-4*B
\end{sagecommandline}



%========================================
\section{Perkalian}

\subsection{Perkalian matriks-vektor}
Perkalian matriks-vektor bekerja seperti yang Anda duga.
\begin{sagecommandline}
sage: A = matrix(QQ, [[1, 3, 5, 9], [0, 2, 4, 6]])
sage: v = vector(QQ, [1, 2, 3, 4])
sage: A, v
sage: A*v
\end{sagecommandline}
Matriks $\nbym{2}{4}$, $A$, mengalikan vektor kolom
$\nbym{4}{1}$, $\vec{v}$, yang diletakkan di sebelah kanan,
seperti pada $A\vec{v}$.
Mencoba menaruh vektor ini di sebelah kiri, seperti di bawah,
menghasilkan
\inlinecode{unsupported operand parent(s) for *}.
\begin{lstlisting}
sage: v*A
\end{lstlisting}

Tentu saja, Anda dapat mengalikan dengan vektor di sebelah kiri
jika ukuran vektor itu sesuai dengan matriksnya.
Di sini, matriks~$A$ yang mempunyai dua baris dapat dikalikan dengan
vektor berentri dua.
\begin{sagecommandline}
sage: w = vector(QQ, [3, 5])
sage: w*A
\end{sagecommandline}




\subsection{Perkalian matriks-matriks}
\Sage{} dapat mengalikan matriks.
\begin{sagecommandline}
sage: A = matrix(QQ, [[2, 1], [4, 3]])
sage: B = matrix(QQ, [[5, 6, 7], [8, 9, 10]]) 
sage: A*B
\end{sagecommandline}
Mencoba $\inlinecode{B*A}$ menghasilkan
\inlinecode{unsupported operand parent(s) for '*'} karena $BA$ tidak terdefinisi.
% \begin{sagecommandline}
% sage: A = matrix(QQ, [[2, 1], [4, 3]])
% sage: B = matrix(QQ, [[5, 6, 7], [8, 9, 10]]) 
% sage: B*A
% \end{sagecommandline}

Untuk matriks persegi berukuran sama, perkalian terdefinisi dalam kedua urutan.
\begin{sagecommandline}
sage: A = matrix(QQ, [[1, 2], [3, 4]])
sage: B = matrix(QQ, [[4, 5], [6, 7]])
sage: A*B
sage: B*A
\end{sagecommandline}
Perhatikan bahwa kedua hasil tersebut berbeda;
perkalian matriks tidak komutatif.
\begin{sagecommandline}[d,0,2]
sage: A*B == B*A
\end{sagecommandline}

Faktanya, perkalian matriks sangat tidak komutatif dalam arti bahwa
jika Anda membuat dua matriks $\nbyn{n}$ secara acak,
keduanya hampir pasti tidak komutatif.
Di sini kita mengalikan seribu pasang matriks acak untuk melihat
apakah ada pasangan yang komutatif.
\begin{sagecommandline}
sage: number_commuting = 0 
sage: for n in range(1000):                                       
....:     A = random_matrix(RDF, 2, min=-1, max=1)
....:     B = random_matrix(RDF, 2, min=-1, max=1)
....:     if (A*B == B*A):
....:         number_commuting = number_commuting + 1 
\end{sagecommandline}
Berikut hasilnya.
\begin{sagecommandline}
sage: number_commuting
\end{sagecommandline}

% Plug a square matrix into a polynomial.
 


\subsection{Balikan}
Jika $A$ adalah matriks taksingular $\nbyn{n}$, maka balikannya $A^{-1}$
adalah matriks $\nbyn{n}$ sedemikian sehingga $A^{-1}A=AA^{-1}$ merupakan
matriks identitas $\nbyn{n}$.
% For $\nbyn{2}$ matrix inverses we have a formula.
Dalam buku, kita memakai rumus untuk mencari balikan matriks $\nbyn{2}$,
sedangkan untuk mencari balikan yang lebih besar,
kita menuliskan matriks asal di sebelah matriks identitas
lalu melakukan reduksi Gauss--Jordan.
\begin{sagecommandline}
sage: A = matrix(QQ, [[1, 3, 1], [2, 1, 0], [4, -1, 0]])
sage: A.is_singular()
sage: I = identity_matrix(3)
sage: B = A.augment(I, subdivide=True)
sage: B
sage: C = B.rref()
sage: C
\end{sagecommandline}
Balikannya berada di sebelah kanan.
Untuk mengambil balikan tersebut,
\Sage{} mempunyai metode \inlinecode{matrix_from_columns}
pada suatu instans matriks yang menerima daftar kolom,
mengekstraknya, lalu mengembalikan matriks yang tersusun dari kolom-kolom itu.
\begin{sagecommandline}
sage: A_inv = C.matrix_from_columns([3, 4, 5])
sage: A_inv
sage: A_inv*A
sage: A*A_inv == identity_matrix(3)
\end{sagecommandline}

Semua langkah ini sangat merepotkan, sehingga tentu tersedia
perintah tersendiri.
\begin{sagecommandline}[d,0,1]
sage: A = matrix(QQ, [[1, 3, 1], [2, 1, 0], [4, -1, 0]])
sage: A_inv = A.inverse()
sage: A_inv
\end{sagecommandline}

Salah satu alasan mencari balikan adalah untuk mempermudah penyelesaian
sistem linear.
Ketiga sistem berikut
\begin{equation*}
  \begin{linsys}{3}
    x  &+ &3y &+ &z &= &4 \\
    2x &+ &y  &  &  &= &4 \\
    4x &- &y  &  &  &= &4 
  \end{linsys}
  \qquad\qquad
  \begin{linsys}{3}
    x  &+ &3y &+ &z &= &2 \\
    2x &+ &y  &  &  &= &-1 \\
    4x &- &y  &  &  &= &5 
  \end{linsys}
  \qquad\qquad
  \begin{linsys}{3}
    x  &+ &3y &+ &z &= &1/2 \\
    2x &+ &y  &  &  &= &0 \\
    4x &- &y  &  &  &= &12 
  \end{linsys}
\end{equation*}
mempunyai matriks koefisien yang sama di ruas kiri,
tetapi ruas kanannya berbeda.
Jika Anda menghitung balikan matriks tersebut,
maka penyelesaian setiap sistem hanya memerlukan satu perkalian matriks-vektor.
\begin{sagecommandline}
sage: A = matrix(QQ, [[1, 3, 1], [2, 1, 0], [4, -1, 0]])
sage: A_inv = A.inverse()
sage: v1 = vector(QQ, [4, 4, 4])
sage: v2 = vector(QQ, [2, -1, 5])
sage: v3 = vector(QQ, [1/2, 0, 12])
sage: A_inv*v1
sage: A_inv*v2
sage: A_inv*v3
\end{sagecommandline}



\section{Bilangan kondisi}
Untuk sistem linear yang dikerjakan dengan tangan dalam buku, kita mencari
solusi eksak.
Namun, dalam penerapan kita menggunakan bilangan titik mengambang sehingga
perhitungan dapat kehilangan presisi akibat persoalan numerik.

Sebagian sistem lebih rentan terhadap kehilangan ini daripada sistem lainnya.
Kerentanan matriks koefisien suatu sistem terhadap persoalan ini
diukur dengan \textit{bilangan kondisi}.
Bilangan ini merupakan bilangan real positif yang
memperkirakan kehilangan presisi dalam kasus terburuk.\footnote{%
  Logaritma basis~$10$ dari bilangan kondisi merupakan perkiraan kasus terburuk
  untuk banyaknya digit yang hilang ketika menyelesaikan sistem linear dengan
  matriks koefisien tersebut.
  Jadi, bilangan kondisi tergolong besar apabila logaritma basis~$10$-nya
  lebih besar daripada atau sama dengan banyaknya digit signifikan
  pada entri-entri matriks.}

Secara lebih tepat, kita mulai dengan sistem linear
$A\vec{v}=\vec{d}$,
lalu menanyakan seberapa besar ketidakakuratan pada~$\vec{d}$ (mungkin akibat
persoalan titik mengambang atau batas pengukuran pada entri-entri vektor)
dapat memengaruhi solusi~$\vec{v}$.
Artinya, kita meninjau
$A(\vec{v}+\vec{\Delta v})=\vec{d}+\vec{\Delta d}$
dan menanyakan hubungan antara $\vec{\Delta d}$ dan $\vec{\Delta v}$.
Linearitas memberikan $A\vec{\Delta v}=\vec{\Delta d}$, sehingga pertanyaan kita
adalah:~seberapa besar $\vec{\Delta v}$ dipengaruhi oleh $\vec{\Delta d}$
melalui matriks~$A$?

Definisikan \textit{norma-$2$} suatu vektor sebagai panjangnya,
  $\norm{\vec{v}}=\sqrt{v_1^2+\cdots+v_n^2}$
(terdapat banyak norma, tetapi inilah yang paling umum).
\begin{sagecommandline}
sage: v = vector(RDF, [1, 2, 3])
sage: v.norm(2)
\end{sagecommandline}
Ada suatu vektor masukan yang direntangkan paling besar oleh $A$, yaitu
vektor $\vec{v}_{\text{max}}$ yang membuat bilangan
$M=\norm{A\vec{v}_{\text{max}}}/\norm{\vec{v}_{\text{max}}}$
maksimal.
Ada pula vektor yang diciutkan paling besar oleh $A$, yaitu
vektor $\vec{v}_{\text{min}}$ yang membuat bilangan
$m=\norm{A\vec{v}_{\text{min}}}/\norm{\vec{v}_{\text{min}}}$ minimal
(dengan anggapan $\vec{v}_{\text{min}}\neq\zero$).\footnote{%
  Faktor-faktor ini akan dibahas lebih lanjut dalam bab berikutnya.}
Bilangan kondisi adalah rasio keduanya,
$\kappa(A)=M/m$, kecuali bahwa matriks singular memetakan beberapa vektor
taknol ke nol sehingga mempunyai $m=0$; dalam hal ini kita menetapkan
$\kappa(A)=\infty$.

Definisi~$M$ memberikan $\norm{\vec{d}}\leq M\norm{\vec{v}}$, sedangkan
definisi~$m$ memberikan $\norm{\vec{\Delta d}}\geq m\norm{\vec{\Delta v}}$.
Dengan demikian, perubahan relatif pada solusi sistem yang dihitung
dan perubahan relatif pada ruas kanan sistem memenuhi perbandingan berikut.
\begin{equation*}
  \frac{\norm{\vec{\Delta v}}}{\norm{\vec{v}}}
   \leq \kappa(A)\cdot \frac{\norm{\vec{\Delta d}}}{\norm{\vec{d}}}
\end{equation*}
Dengan kata lain, bilangan kondisi adalah faktor kasus terburuk yang
memperbesar galat relatif.
Perubahan pada ruas kanan dapat menimbulkan perubahan pada solusi yang dihitung
hingga $\kappa(A)$ kali lebih besar.
\begin{sagecommandline}
sage: A = matrix(RDF, 3, [[1, 2, 3], [4, 5, 6], [7, 8, 9.001]])
sage: A.condition(2)
\end{sagecommandline}
Matriks yang hampir singular mempunyai bilangan kondisi lebih besar daripada
matriks yang jauh dari singular.
\begin{sagecommandline}
sage: v = []                                                                    
sage: for i in range(300): 
....:     A = random_matrix(RDF, 3, min=0, max=100) 
....:     v.append(A.condition(2))                                                                           
\end{sagecommandline}
Kini kita dapat memperoleh gambaran tentang nilai yang lazim.
\begin{sagecommandline}
sage: mean(v), std(v)
\end{sagecommandline}
Jika bilangan kondisinya besar, kita mengatakan bahwa
sistem untuk persoalan tersebut `berkondisi buruk';
jika tidak, kita mengatakan bahwa sistem itu `berkondisi baik'.





\section{Waktu eksekusi}
Persoalan aljabar linear berukuran besar sering dijumpai dalam sains dan
rekayasa.
Karena komputer bekerja cepat dan akurat,
komputer memungkinkan kita menyelesaikan persoalan
yang mustahil kita kerjakan dengan tangan.

Namun, komputer pun mempunyai keterbatasan.
Salah satu batas ukuran persoalan yang dapat kita tangani ditentukan oleh
seberapa cepat mesin dapat menemukan jawabannya.
Wajarnya, persoalan dengan matriks lebih besar cenderung memerlukan waktu
lebih lama untuk diselesaikan.
Membandingkan ukuran persoalan dengan waktu yang diperlukan suatu algoritma
untuk menyelesaikannya merupakan cara penting untuk menilai kegunaan
algoritma tersebut.

Operasi balikan matriks merupakan ilustrasi yang baik.
% This is an important operation; for instance, if we could do large matrix 
% inverses
% quickly then as show above we could quickly solve large linear systems.
(Entri-entri matriks ini berupa bilangan titik mengambang
karena jenis inilah yang paling umum dalam penerapan.)
% Timing gives bad output.  It has extra mu's in there.  I don't
% know how to fix it; I suspect it is a bug in sagemath.sty or listings.sty's
% inability to do non-ASCII.
% \begin{sagecommandline}
% sage: A = matrix(RDF, [[1, 3, 1], [2, 1, 0], [4, -1, 0]])
% sage: A
% sage: A.is_singular()
% sage: timeit('A.inverse()')
% \end{sagecommandline}
\begin{lstlisting}
sage: A = matrix(RDF, [[1, 3, 1], [2, 1, 0], [4, -1, 0]])
sage: A
[ 1.0  3.0  1.0]
[ 2.0  1.0  0.0]
[ 4.0 -1.0  0.0]
sage: A.is_singular()
False
sage: timeit('A.inverse()')
625 loops, best of 3: 62.7 micro-s per loop
\end{lstlisting}
% was: 625 loops, best of 3: 62.7 μs per loop
Dalam \Sage, \inlinecode{timeit} memberi tahu berapa lama
suatu operasi berlangsung.
Perintah ini tidak hanya melaporkan waktu yang dipakai CPU,
melainkan lama operasi sebagaimana diukur oleh jam sebenarnya.
Perintah tersebut menjalankan operasi berkali-kali untuk mengurangi
dampak perlambatan komputer akibat penulisan ke cakram atau gangguan lain.
Baris terakhir menyatakan bahwa perintah dijalankan dalam tiga kelompok,
masing-masing sebanyak $625$ kali, dan waktu rata-rata dalam kelompok
tercepat adalah $62.7$~mikrodetik.\footnote{% from cat /proc/cpuinfo
  Mesin ini mengidentifikasi dirinya sebagai:
  Intel(R) Core(TM) i7-6820HQ CPU @ 2.70GHz.}
Itu cepat, tetapi $A$ memang hanya berukuran $\nbyn{3}$.

Terlebih lagi, $A$ adalah satu matriks $\nbyn{3}$ tertentu.
Untuk memperkirakan waktu yang biasanya diperlukan,
Anda dapat mencoba mencari balikan matriks acak.
% \begin{sagecommandline}
% sage: timeit('random_matrix(RDF, 3, min=-1, max=1).inverse()')
% sage: timeit('random_matrix(RDF, 3, min=-1, max=1).inverse()')
% sage: timeit('random_matrix(RDF, 3, min=-1, max=1).inverse()')
% \end{sagecommandline}
\begin{lstlisting}
sage: timeit('random_matrix(RDF, 3, min=-1, max=1).inverse()')
625 loops, best of 3: 86 micro-s per loop
sage: timeit('random_matrix(RDF, 3, min=-1, max=1).inverse()')
625 loops, best of 3: 85.8 micro-s per loop
sage: timeit('random_matrix(RDF, 3, min=-1, max=1).inverse()')
625 loops, best of 3: 86.6 micro-s per loop
\end{lstlisting}
% was: 
%% sage: timeit('random_matrix(RDF, 3, min=-1, max=1).inverse()')
%% 625 loops, best of 3: 86 μs per loop
%% sage: timeit('random_matrix(RDF, 3, min=-1, max=1).inverse()')
%% 625 loops, best of 3: 85.8 μs per loop
%% sage: timeit('random_matrix(RDF, 3, min=-1, max=1).inverse()')
%% 625 loops, best of 3: 86.6 μs per loop

Salah satu masalah pada data kinerja ini adalah bahwa
kita tidak dapat mengetahui berapa banyak waktu yang dipakai untuk membuat
matriks acak dan berapa banyak yang dipakai untuk mencari balikannya.
Kode di bawah mengambil beberapa ukuran, $\nbyn{3}$, $\nbyn{10}$, dan seterusnya,
membuat satu matriks acak, lalu mengukur waktu untuk menghitung
balikan matriks tersebut.
\begin{lstlisting}
sage: for size in [3, 10, 25, 50, 75, 100, 150, 200]:
....:     print("ukuran = "+str(size))
....:     M = random_matrix(RR, size, min=-1, max=1)
....:     timeit('M.inverse()')
....:     
ukuran = 3
625 loops, best of 3: 55 micro-s per loop
ukuran = 10
625 loops, best of 3: 548 micro-s per loop
ukuran = 25
125 loops, best of 3: 6.79 ms per loop
ukuran = 50
5 loops, best of 3: 53 ms per loop
ukuran = 75
5 loops, best of 3: 177 ms per loop
ukuran = 100
5 loops, best of 3: 419 ms per loop
ukuran = 150
5 loops, best of 3: 1.41 s per loop
ukuran = 200
5 loops, best of 3: 3.4 s per loop
\end{lstlisting}
% was:
%% size = 3
%% 625 loops, best of 3: 55 μs per loop
%% size = 10
%% 625 loops, best of 3: 548 μs per loop
%% size = 25
%% 125 loops, best of 3: 6.79 ms per loop
%% size = 50
%% 5 loops, best of 3: 53 ms per loop
%% size = 75
%% 5 loops, best of 3: 177 ms per loop
%% size = 100
%% 5 loops, best of 3: 419 ms per loop
%% size = 150
%% 5 loops, best of 3: 1.41 s per loop
%% size = 200
%% 5 loops, best of 3: 3.4 s per loop
Sebagian waktu tersebut dinyatakan dalam mikrodetik, sebagian dalam
milidetik, dan sebagian dalam detik.
Satu mikrodetik adalah sepersejuta detik,
$0.000\,001$~detik.
Satu milidetik adalah seperseribu detik,
$0.001$~detik.
Tabel berikut menyatakan semuanya dalam detik.
\begin{center}
  \begin{tabular}{r|r@{.}l}
    \multicolumn{1}{r}{\textit{ukuran}}     &\multicolumn{2}{c}{\textit{detik}}  \\  \hline
    $3$      &$0$ &$000\,055$ \\
    $10$     &$0$ &$000\,548$ \\
    $25$     &$0$ &$006\,79$ \\
    $50$     &$0$ &$053$ \\
    $75$     &$0$ &$177$ \\
    $100$    &$0$ &$419$ \\
    $150$    &$1$ &$41$ \\
    $200$    &$3$ &$4$ 
  \end{tabular}
\end{center}
Waktu bertambah lebih cepat daripada ukuran.
Misalnya, menggandakan ukuran dari~$25$ menjadi~$50$ meningkatkan waktu
lebih dari dua kali lipat: $0.053/0.00679$ kira-kira $7.8$.
Demikian pula, memperbesar ukuran empat kali lipat dari $50$ menjadi~$200$
membuat waktu meningkat jauh lebih dari empat kali lipat:
$3.4/0.053\approx 64.15$.
Perubahan ukuran persoalan dari~$10$ menjadi~$100$ meningkatkan waktu yang
diperlukan dengan faktor $0.419/0.000\,548\approx 764.60$.

Peroleh grafik dengan memberikan data kepada \Sage{} sebagai daftar pasangan.\footnote{% 
  Grafik dalam panduan ini dibuat dengan menggunakan lebih banyak opsi
  penggambaran daripada yang tampak pada blok keluaran.
  Misalnya, diagram pencar di sini dibuat dengan
  \protect\inlinecode{g = scatter_plot(d, markersize=10, facecolor='#b9b9ff')}
  dan disimpan ke dalam berkas dengan
  \protect\inlinecode{g.save("graphics/mat001.pdf", figsize=[2.25,1.5], axes_pad=0.05, fontsize=7)}.
  Sebagian besar kode penghias ini akan kita hilangkan agar tidak mengganggu.}
\begin{sagecommandline}
sage: d = [(3, 0.000055), (10, 0.000548), (25, 0.00679),  
....:      (50, 0.053), (75, 0.177), (100, 0.419), 
....:      (150, 1.41), (200, 3.4)]
sage: g = scatter_plot(d)  
sage: g.save("graphics/mat001.pdf")            
\end{sagecommandline}
\begin{sagesilent}
g = scatter_plot(d, markersize=10, facecolor='#b9b9ff')
g.save("graphics/mat001.pdf", figsize=[2.25,1.5], axes_pad=0.05, fontsize=7)
\end{sagesilent}
(Jika Anda memasukkan \inlinecode{scatter_plot(d)} pada prompt, yaitu
tanpa menyimpannya sebagai~$g$, maka \Sage{} akan membuka jendela yang
menampilkan grafik.)
\begin{center}
  \includegraphics{graphics/mat001.pdf}
\end{center}
Grafik tersebut memperjelas bahwa rasio $\text{waktu}/\text{ukuran}$
tidak konstan karena data jelas tidak terletak pada satu garis lurus.

Berikut beberapa data tambahan.
Perlu diingat jika Anda mencobanya sendiri:~untuk menghasilkan data ini,
\inlinecode{timeit} memerlukan waktu begitu lama sehingga komputer harus
dibiarkan berjalan semalaman.
\begin{lstlisting}
sage: for size in [500, 750, 1000]:                             
....:     print("ukuran=", size)
....:     M = random_matrix(RR, size, min=-1, max=1)
....:     timeit('M.inverse()')
....: 
ukuran= 500
5 loops, best of 3: 51.4 s per loop
ukuran= 750
5 loops, best of 3: 172 s per loop
ukuran= 1000
5 loops, best of 3: 406 s per loop
\end{lstlisting}
Sekali lagi, tabel merupakan cara yang lebih rapi untuk menyajikan data.
\begin{center}
  \begin{tabular}{r|r@{.}l}
    \multicolumn{1}{r}{\textit{ukuran}}     &\multicolumn{2}{c}{\textit{detik}}  \\  \hline
    $500$       &$51$ &$4$ \\
    $750$       &$172$ &   \\
    $1000$      &$406$ &   
  \end{tabular}
\end{center}
Buat grafik dengan menambahkan data baru ke data yang sudah ada.
\begin{sagecommandline}
sage: d = [(3, 0.000055), (10, 0.000548), (25, 0.00679),  
....:      (50, 0.053), (75, 0.177), (100, 0.419), 
....:      (150, 1.41), (200, 3.4)]
sage: d = d + [(500, 51.4), (750, 172), (1000, 406)]
sage: g = scatter_plot(d)                           
sage: g.save("graphics/mat002.pdf")                      
\end{sagecommandline}
\begin{sagesilent}
# d = [(3, 0.000055),    
#      (10, 0.000548), 
#      (25, 0.00679),  
#      (50, 0.053), 
#      (75, 0.177), 
#      (100, 0.419), 
#      (150, 1.41), 
#      (200, 3.4)]
# d = d + [(500, 51.4), (750, 172), (1000, 406)]
g = scatter_plot(d, markersize=10, facecolor='#b9b9ff')
g.save("graphics/mat002.pdf", figsize=[2.25,2.25], axes_pad=0.05, fontsize=7)              
\end{sagesilent}
% The result is this image.
\begin{center}
  \includegraphics{graphics/mat002.pdf}
\end{center}
(Perhatikan bahwa kedua grafik mempunyai skala yang berbeda;
jika grafik pada halaman ini memakai skala vertikal yang sama dengan grafik
pada halaman sebelumnya, grafik ini akan menjulur jauh melewati bagian atas
kertas.)
Jadi, batas praktis ukuran persoalan yang dapat kita selesaikan dengan
operasi balikan matriks timbul dari kenyataan bahwa grafik di atas bukan
garis lurus\Dash waktu yang diperlukan bertambah jauh lebih cepat daripada
ukurannya.

Salah satu upaya utama dalam Ilmu Komputer adalah menemukan algoritma cepat
untuk mengerjakan tugas praktis.
Banyak orang secara khusus mengerjakan persoalan Aljabar Linear,
seperti mencari balikan matriks, karena persoalan semacam itu sangat umum
dalam penerapan.

\endinput


TODO:
